% ch05-ai-concepts
\section{الذكاء الاصطناعيّ (1): المفاهيم الأساسية}

في هذا الفصل نبني الأساس النظريّ الذي يقوم عليه كلّ ما يخصّ الذكاء الاصطناعيّ في
\en{ScholarPath}. لن نكتب كودًا كثيرًا هنا؛ هدفنا أن نفهم \emph{الأفكار} أولًا:
كيف يحوّل الحاسوب النصّ إلى أرقام؟ وكيف يقيس أنّ جملتين متقاربتان في المعنى؟ وما
\en{RAG} ولماذا نحتاجه أصلًا؟ ثمّ نُنهي الفصل بالصورة الكاملة: عندنا ثلاث ميزات
ذكاء اصطناعيّ، لكنّ \textbf{واحدة فقط} منها تستعمل نموذجًا لُغويًّا حقيقيًّا. نمشي
واحدة واحدة، وكلّ فكرة نربطها بسطر كود حقيقيّ من المشروع.

% ─────────────────────────────────────────────────────────────────────────────
\subsection{الـ Embeddings — تحويل النصّ إلى أرقام}

\begin{idea}[الفكرة: \en{Embedding}]
الـ \en{Embedding} يأخذ أيّ نصّ — كلمة أو جملة أو فقرة — ويحوّله إلى مصفوفة أرقام
عشريّة ثابتة الطول اسمها \en{vector} (مُتّجه). الخاصّيّة الجوهريّة: \textbf{النصوص
المتقاربة في المعنى تُنتج مُتّجهات متقاربة عدديًّا}، بغضّ النظر عن الصياغة الحرفيّة،
بل وحتى عبر اللغات. سؤال بالإنجليزيّة ونفس السؤال بالعربيّة يُنتجان مُتّجهين يشيران
إلى نفس الاتّجاه تقريبًا داخل فضاء ذي 1536 بُعدًا.
\end{idea}

تخيّل خريطة ضخمة، لكن بدل بُعدين (طول وعرض) فيها 1536 بُعدًا. كلّ نصّ يصير نقطةً على
هذه الخريطة. المنح الدراسيّة تتجمّع في ناحية، وأسئلة الأمور الشائعة في ناحية، وملفّات
المستشارين في ناحية ثالثة. لمّا يسأل الطالب سؤالًا، نحوّل سؤاله إلى نقطة أيضًا، ثمّ
نبحث عن أقرب النقاط إليه — وتلك النقاط هي المستندات الأوثق صِلةً بسؤاله.

في \en{ScholarPath} نموذج الـ \en{Embeddings} هو \en{Azure OpenAI}
\code{text-embedding-3-small}، وهو يُنتج مُتّجهات بطول \textbf{1536} رقمًا. لكن هنا
تفصيلة مهمّة: لو خزّنّا المُتّجهات في القاعدة دون أن نُسجّل \emph{أيّ نموذج} ولّدها،
فسنقع في كارثة صامتة يوم نغيّر النموذج — لأنّ مُتّجهات نموذجٍ لا تُقارَن أبدًا
بمُتّجهات نموذجٍ آخر. الحلّ بسيط: نضع بادئة \code{sp:} أمام اسم النموذج، فتصير هويّة
المُتّجه واضحة ومحجوزة.

\textbf{هنشوف إزّاي الكود بيبني اسم النموذج بالبادئة \code{sp:}} — المقتطف من الملفّ
\code{AzureOpenAiEmbeddingService.cs}:

\begin{english}
\begin{lstlisting}[language={[Sharp]C}]
// Provider-neutral model tag. The prefix namespaces the vector space so
// documents embedded with a different model are never matched here.
public string ModelName =>
    IsAzureEmbeddingConfigured()
        ? $"sp:{opts.Value.AzureOpenAi.EmbeddingDeploymentName}"
        // -> "sp:text-embedding-3-small"
        : local.ModelName;   // fallback when Azure isn't configured

public int Dimensions => opts.Value.AzureOpenAi.EmbeddingDimensions; // 1536
\end{lstlisting}
\end{english}

نقرأ الكود سطرًا سطرًا:
\begin{itemize}
  \item \code{ModelName} خاصّيّة محسوبة: لو \en{Azure} مُهيّأ، تُرجِع الاسم مسبوقًا
    بـ \code{sp:} (فيصير \code{sp:text-embedding-3-small}).
  \item لو \en{Azure} غير مُهيّأ، ترجع لاسم النموذج المحلّيّ عبر \code{local.ModelName}
    — هذا مبدأ الاختيار بالإعداد الذي عرفناه سابقًا.
  \item \code{Dimensions} تُرجِع 1536، وهو طول المُتّجه الذي نلتزم به في القاعدة.
\end{itemize}

\begin{note}
البادئة \code{sp:} ليست جزءًا من اسم النموذج عند \en{Azure}؛ هي مجرّد «مساحة أسماء»
(\en{namespace}) داخل قاعدتنا. فائدتها الوحيدة: تمييز أيّ نموذجٍ ولّد كلّ مُتّجه،
حتى لا نخلط مُتّجهات نموذجين مختلفين في بحثٍ واحد.
\end{note}

% ─────────────────────────────────────────────────────────────────────────────
\subsection{الـ Cosine Similarity — قياس التقارب المعنويّ}

عرفنا أنّ النصوص المتقاربة تُنتج مُتّجهات متقاربة الاتّجاه. لكن كيف نقيس هذا التقارب
بالضبط؟ الجواب: نقيس \emph{الزاوية} بين المُتّجهين، لا المسافة بينهما ولا طولهما.
هذا هو \en{Cosine Similarity} (تشابه جيب التمام). لو كان لدينا مُتّجهان
$\mathbf{a}$ و $\mathbf{b}$، فالمعادلة:

\[
\cos(\theta)
= \frac{\mathbf{a}\cdot\mathbf{b}}{\|\mathbf{a}\|\;\|\mathbf{b}\|}
= \frac{\displaystyle\sum_{i=1}^{n} a_i\, b_i}
       {\displaystyle\sqrt{\sum_{i=1}^{n} a_i^2}\;\sqrt{\sum_{i=1}^{n} b_i^2}}
\quad\in[0,\,1]
\]

\begin{idea}[الفكرة: الزاوية الصغيرة = تشابه عالٍ]
البَسط $\mathbf{a}\cdot\mathbf{b}$ هو الضرب النقطيّ، والمقام يقسم على طولَي
المُتّجهين حتى يُلغي أثر الطول تمامًا ويُبقي الاتّجاه وحده. النتيجة رقم بين $0$
و $1$: كلّما \textbf{صغُرت الزاوية} بين المُتّجهين اقتربت النتيجة من $1$ (تشابه
عالٍ)، وكلّما كبُرت الزاوية اقتربت من $0$ (لا علاقة). ببساطة: مُتّجهان يشيران لنفس
الاتّجاه $\Rightarrow$ نصّان يتحدّثان عن نفس الموضوع.
\end{idea}

لكن ليس كلّ تقاربٍ يُقبَل. نحتاج \emph{عتبة} (\en{threshold}) نقول عندها: «هذا
المستند قريب بما يكفي لِنُرفقه للنموذج». هذه العتبة مضبوطة في الإعداد
\code{RagMinScore} وقيمتها \textbf{0.6}. الجدول التالي يوضّح كيف نُفسّر الدرجات:

\begin{center}
\begin{tabular}{L{2.8cm} L{4.2cm} L{6.6cm}}
\toprule
\textbf{\color{spPrimary}الدرجة} & \textbf{\color{spPrimary}التفسير} & \textbf{\color{spPrimary}السلوك في \en{ScholarPath}}\\
\midrule
$\geq 0.85$ & معنى شبه مطابق & يُسترجَع ويُستشهَد به في الإجابة\\
$0.70$–$0.84$ & متشابه جدًّا & يُسترجَع ويُستشهَد به\\
$0.60$–$0.69$ & مرتبط بوضوح & يُسترجَع (عند حدّ \code{RagMinScore})\\
$< 0.60$ & ارتباط ضعيف & يُستبعَد قبل إرسال البرومبت للنموذج\\
\bottomrule
\end{tabular}
\end{center}

\begin{keyp}[نقطة مهمة: العتبة تحمي الجودة]
لو أرسلنا للنموذج مستندات ضعيفة الصِّلة، سيبني إجابته على معلومات لا تخصّ السؤال —
وهذا مصدرٌ رئيسيّ للهلوسة. العتبة \code{RagMinScore = 0.6} هي صمّام الأمان: أيّ مستند
تحت 0.6 يُرمى قبل أن يرى النموذجُ سؤالَ المستخدم أصلًا.
\end{keyp}

% ─────────────────────────────────────────────────────────────────────────────
\subsection{ما هو الـ RAG؟}

\begin{idea}[الفكرة: \en{Retrieval-Augmented Generation}]
النموذج اللُّغويّ (\en{LLM}) يعرف الكثير، لكن معرفته مُجمّدة عند لحظة تدريبه، ولا
يعرف شيئًا عن بيانات \en{ScholarPath} الخاصّة (منحنا، مستشارونا، مقالاتنا). لو سألناه
مباشرةً «ما المنح المتاحة لمصر؟» فسيُخمّن من ذاكرته — وقد يخترع منحًا لا وجود لها.
الـ \en{RAG} يحلّ هذه \emph{فجوة المعرفة}: بدل التخمين، نبحث أوّلًا في قاعدة معرفتنا
الحيّة عن المستندات الأوثق صِلةً بالسؤال، ثمّ نُرفقها للنموذج \textbf{كسياق} مع
السؤال، فتُبنى الإجابة على بياناتٍ حقيقيّة ومحدَّثة.
\end{idea}

الاسم نفسه يشرح المكوّنات الثلاثة: \en{Retrieval} (استرجاع المستندات)، ثمّ
\en{Augmented} (تعزيز البرومبت بها)، ثمّ \en{Generation} (توليد الإجابة). التدفّق في
\en{ScholarPath} \textbf{خطّيّ ومرّة واحدة} — لا حلقات ولا تكرار:

\begin{center}
\begin{tabular}{L{1.4cm} L{12.4cm}}
\toprule
\textbf{\color{spPrimary}الخطوة} & \textbf{\color{spPrimary}ما يحدث}\\
\midrule
1 & يكتب المستخدم سؤالًا نصّيًّا.\\
2 & نحوّل السؤال إلى مُتّجه (\code{EmbedAsync} $\Rightarrow$ \code{float[1536]}).\\
3 & نحسب \en{cosine} بين مُتّجه السؤال وكلّ مستند، ونأخذ الأعلى تشابهًا.\\
4 & نُبقي المستندات فوق العتبة (0.6) فقط، ونصوغها في كتلة سياق.\\
5 & نُرسل: تعليمات النظام + السياق + السؤال إلى النموذج \en{GPT-4o-mini}.\\
6 & يُرجِع النموذج إجابةً مسنودةً، مع قائمة المصادر التي استُشهد بها.\\
\bottomrule
\end{tabular}
\end{center}

\begin{note}
لاحظ الفرق الجوهريّ: النموذج لا «يتذكّر» المنح من تدريبه، بل \emph{يقرؤها} من السياق
الذي حضّرناه له لحظتها. لو أضاف الأدمِن منحةً جديدة اليوم وأعاد بناء قاعدة المعرفة،
صار النموذج قادرًا على ذكرها فورًا — دون إعادة تدريبٍ ولا تعديل كود.
\end{note}

% ─────────────────────────────────────────────────────────────────────────────
\subsection{الـ Standard RAG مقابل الـ Agentic RAG}

التدفّق الذي شرحناه يُسمّى \en{Standard RAG}: خطّ أنابيب مستقيم، تمريرة واحدة، بلا
خطوة تقييم. هذا بسيط وسريع، لكن له ثلاثة أنماط فشل معروفة:

\begin{enumerate}
  \item \textbf{الأسئلة الغامضة}: «إزّاي أقدّم؟» سؤال ناقص السياق — هل يقصد التسجيل في
    المنصّة أم التقديم على منحة بعينها؟ الـ \en{Standard RAG} يسترجع دون أن يستوضح.
  \item \textbf{الأدلّة المبعثرة}: سؤال يحتاج تجميعًا من عدّة عناقيد مستندات (مثل «قارِن
    \en{Erasmus} بـ \en{Chevening} في التمويل») يحصل على إجابةٍ جزئيّة، لأنّ خطّ
    الأنابيب يستنفد تمريرة استرجاعٍ واحدة فقط.
  \item \textbf{الثقة الزائفة}: مستندات سجّلت تشابهًا عاليًا لكنّها لا تحوي الإجابة فعلًا،
    فيملأ النموذج الفراغ بتفاصيل مُقنعة لكنّها مُختلَقة.
\end{enumerate}

الـ \en{Agentic RAG} يعالج هذا باستبدال الخطّ المستقيم بـ \textbf{حلقة تحكّم}
(\en{control loop}) يقودها وكيلٌ (\en{agent}) ذاتيّ القرار: يستوضح لو كان السؤال
غامضًا، ويختار مصدر البحث المناسب، ويُقيّم هل النتائج كافية، ولو لم تكن كافية يُعيد
صياغة السؤال ويبحث مجدّدًا — قبل أن يُولّد الإجابة النهائيّة. أي أنّه يكرّر ويصحّح
نفسه، لكن بثمن. الجدول التالي يوازن بينهما:

\begin{center}
\begin{tabularx}{\linewidth}{L{3.6cm} Y Y}
\toprule
\textbf{\color{spPrimary}المعيار} & \textbf{\color{spPrimary}\en{Standard RAG}} & \textbf{\color{spPrimary}\en{Agentic RAG}}\\
\midrule
زمن الاستجابة & 1–3 ثوانٍ & 5–15+ ثانية (نداءات نموذج متعدّدة)\\
التكلفة لكلّ استعلام & $\approx \$0.0002$ & $\approx \$0.001$–$\$0.002$ ($3$–$10$ أضعاف)\\
سؤال بسيط أحاديّ المصدر & ممتاز & بلا تحسّن يُذكَر\\
سؤال معقّد متعدّد المصادر & متوسّط & أفضل بوضوح\\
تصحيح الأخطاء & حتميّ وقابل للاختبار & غير حتميّ، يصعُب تكراره\\
حجم المتون الأمثل & حتى $\sim$10 آلاف مستند & 10 آلاف+ ومجالات متعدّدة\\
\bottomrule
\end{tabularx}
\end{center}

\begin{why}
اخترنا \en{Standard RAG} للشات بوت. قاعدة معرفتنا فيها \textbf{692} مستندًا فقط،
والأسئلة موضوعيّة ومحدّدة (منح، مستشارون، نصائح تقديم). تمريرة استرجاعٍ واحدة تجد
المستند الصحيح باستمرار في \textbf{أقلّ من 10 ملّي ثانية}. إضافة حلقة \en{agentic}
كانت ستضرب الزمن في $5$ والتكلفة في $10$ دون أيّ مكسبٍ ملموس في الدقّة على مجموعة
أسئلتنا. عتبة الانتقال إلى \en{Agentic RAG} تقارب 10 آلاف مستندٍ بمجالاتٍ متغايرة —
ونحن بعيدون عنها.
\end{why}

\begin{qa}
\textbf{س: أليس \en{Agentic RAG} «أذكى» دائمًا؟ لماذا لا نستعمله لِنبدو متطوّرين؟}\\
لأنّ الهندسة الجيّدة تختار الأبسط الذي يفي بالغرض. «الأذكى» هنا يعني أبطأ بخمسة
أضعاف، وأغلى بعشرة، وأصعب في الاختبار وتتبّع الأخطاء — مقابل صفر تحسّنٍ قابلٍ للقياس
على 692 مستندًا. القرار الهندسيّ السليم أن نضيف التعقيد يوم تفرضه البيانات، لا قبله.
\end{qa}

% ─────────────────────────────────────────────────────────────────────────────
\subsection{الصورة الكاملة — ثلاث ميزات، مسارٌ واحد فقط يستعمل GPT}

\begin{keyp}[نقطة مهمة: تمييزٌ جوهريّ]
في \en{ScholarPath} \textbf{ثلاث} ميزات ذكاء اصطناعيّ، لكن \textbf{واحدة فقط}
تستعمل \en{RAG} + \en{GPT} حقيقيًّا (وهي الشات بوت). الميزتان الأخريان — التوصيات
والأهليّة — حسابيّتان بالكامل: خوارزميّة تسجيل نقاط حتميّة تعمل محلّيًّا داخل الخادم،
بلا أيّ نداء شبكة ولا نموذجٍ لُغويّ.
\end{keyp}

الجدول التالي يلخّص التوجيه (\en{routing}):

\begin{center}
\begin{tabular}{L{3.4cm} C{3.4cm} C{2.6cm} L{4.0cm}}
\toprule
\textbf{\color{spPrimary}الميزة} & \textbf{\color{spPrimary}تستعمل \en{GPT}؟} & \textbf{\color{spPrimary}تستعمل \en{RAG}؟} & \textbf{\color{spPrimary}الأسلوب}\\
\midrule
الشات بوت & نعم — \en{GPT-4o-mini} & نعم (بحث \en{cosine}) & \en{Standard RAG}\\
التوصيات & لا & لا & تسجيل نقاط محلّيّ\\
الأهليّة & لا & لا & مطابقة ملفٍّ محلّيّة\\
\bottomrule
\end{tabular}
\end{center}

كيف يظهر هذا في الكود؟ الخدمة \code{AzureOpenAiService} هي الواجهة الرسميّة حين يكون
\code{Ai:Provider = AzureOpenAi}، لكنّها لا تُنفّذ التوصيات والأهليّة بنفسها — بل
\textbf{تفوّضها} فورًا إلى \code{LocalAiService}. فقط \code{AskAsync} (الشات) هي التي
تشغّل مسار \en{RAG} الحقيقيّ.

\textbf{هنشوف التفويض في الكود} — المقتطف من \code{AzureOpenAiService.cs}:

\begin{english}
\begin{lstlisting}[language={[Sharp]C}]
// Recommendations & eligibility are structured profile-matching tasks:
// deterministic scoring beats LLM guesswork, so delegate to the local service.
public Task<AiRecommendationResult> GenerateRecommendationsAsync(
    Guid userId, int topN, CancellationToken ct)
    => local.GenerateRecommendationsAsync(userId, topN, ct);

public Task<AiEligibilityResult> CheckEligibilityAsync(
    Guid userId, Guid scholarshipId, CancellationToken ct)
    => local.CheckEligibilityAsync(userId, scholarshipId, ct);

// Only AskAsync runs the real RAG pipeline (retrieve -> augment -> generate).
public async Task<AiChatResponse> AskAsync(...) { /* ... */ }
\end{lstlisting}
\end{english}

نقرأ المقتطف:
\begin{itemize}
  \item الدالّتان \code{GenerateRecommendationsAsync} و \code{CheckEligibilityAsync}
    ليستا إلا سطر تمرير: \code{=> local.\dots} — أي أنّهما تُعيدان النداء حرفيًّا إلى
    \code{LocalAiService} بلا أيّ نداء \en{Azure}.
  \item الدالّة \code{AskAsync} وحدها هي التي تحمل جسدًا حقيقيًّا (\code{async \dots})
    فيه خطوات \en{RAG}. هذا هو الفرق العمليّ بين «ميزة تستعمل \en{GPT}» وأخرى لا تستعمله.
\end{itemize}

والآن ننظر داخل \code{AskAsync} نفسها لِنرى خطوتَي «الاسترجاع» و«الترشيح» — وهنا
يظهر الإعدادان اللذان تحدّثنا عنهما، \code{RagTopK} و \code{RagMinScore}:

\begin{english}
\begin{lstlisting}[language={[Sharp]C}]
// Retrieve (the "R" in RAG): the top-K most similar knowledge docs.
docs = await retriever.RetrieveAsync(msg, ai.RagTopK, ct);   // RagTopK = 5

// Keep only docs above the similarity threshold before building the prompt.
var relevant = docs.Where(d => d.Score >= ai.RagMinScore).ToList(); // 0.6
var context  = RagSupport.BuildContextBlock(relevant, arabic);
\end{lstlisting}
\end{english}

نقرأ:
\begin{itemize}
  \item \code{RetrieveAsync(msg, ai.RagTopK, ct)} تُرجِع أعلى \code{RagTopK} مستندًا
    تشابهًا مع السؤال — والقيمة \textbf{5}. أي نأخذ أفضل خمسة مرشّحين لا أكثر.
  \item \code{Where(d => d.Score >= ai.RagMinScore)} تُرشّح المرشّحين بحدّ \textbf{0.6}؛
    فلو لم يبلغ أيٌّ منها العتبة، تخرج كتلة سياقٍ فارغة، ويُجيب النموذج بحذرٍ بدل أن
    يخترع.
  \item \code{BuildContextBlock} تصوغ الناجين في نصٍّ مُرقّم يُلصَق داخل البرومبت.
\end{itemize}

\begin{secnote}
تفويض التوصيات والأهليّة محلّيًّا ليس اختصارًا فقط، بل قرار خصوصيّة وأمان: بيانات
الطالب الحسّاسة (مستواه الأكاديميّ، بلده، تخصّصه) تُطابَق داخل خادمنا ولا تغادره إلى
أيّ خدمة خارجيّة. لا يخرج إلى \en{Azure} إلا نصّ سؤال الشات وسياقٌ عامّ من قاعدة
المعرفة — لا ملفّات المستخدمين.
\end{secnote}

\begin{summary}[الخلاصة في سطور]
\begin{itemize}
  \item الـ \en{Embedding} يحوّل النصّ إلى مُتّجه ثابت الطول (1536 رقمًا عبر
    \code{text-embedding-3-small})؛ النصوص المتقاربة معنى تُنتج مُتّجهات متقاربة
    اتّجاهًا، حتى عبر اللغات.
  \item بادئة \code{sp:} تحجز «هويّة النموذج» في القاعدة حتى لا نخلط مُتّجهات نموذجين.
  \item الـ \en{Cosine Similarity} يقيس الزاوية بين مُتّجهين (رقم بين 0 و1)؛ الزاوية
    الصغيرة تعني تشابهًا عاليًا، والعتبة \code{RagMinScore = 0.6} تُبعد الضعيف قبل
    النموذج.
  \item الـ \en{RAG} يسدّ فجوة معرفة النموذج: نسترجع سياقًا حيًّا ونُرفقه بالسؤال بدل
    التخمين من الذاكرة — تدفّق خطّيّ واحد.
  \item اخترنا \en{Standard RAG} لا \en{Agentic}: 692 مستندًا واسترجاع أقلّ من 10
    ملّي ثانية لا يبرّران مضاعفة الزمن $5\times$ والتكلفة $10\times$.
  \item ثلاث ميزات \en{AI}، لكن الشات بوت وحده يستعمل \en{RAG}+\en{GPT}؛ التوصيات
    والأهليّة حسابيّتان محلّيًّا عبر \code{LocalAiService}.
\end{itemize}
\end{summary}
